\chapter{Discussion}
\label{chap:discussion}

%% PURPOSE: This chapter interprets the results from Chapter 5 in light of the
%% research questions, prior literature (Chapters 2-3), and the hypotheses
%% (Chapter 1). It also honestly acknowledges what the study cannot claim.
%% The discussion should be the most intellectually dense chapter — avoid
%% re-stating results, focus on what they mean.

\section{Interpretation of Findings}
\label{sec:interpretation}
% TODO: Synthesize the five sets of results into a coherent narrative about
% what virtual choreographies reveal about K-12 online learning behavior.
% Key questions to address:
%   - Do the four discovered archetypes (Steady_Worker, Balanced_Engager,
%     Low_Engagement, Night_Crammer) correspond to theoretically meaningful
%     behavioral types, or are they data artifacts?
%   - How do the choreographies align with constructs from the literature, such
%     as self-regulated learning (SRL) theory and transactional distance?
%   - What does the moderate ARI (~0.53) tell us about the inherent difficulty
%     of recovering behavioral types from noisy, week-level log data?
%   - What does the high temporal stability (week-over-week ARI) imply about
%     whether choreographies are habits or situational responses?

\section{Revisiting the Hypotheses (H1--H5)}
\label{sec:hypotheses-revisited}
% TODO: For each hypothesis, state the verdict (Supported / Partially Supported /
% Not Supported / Inconclusive) and cite the specific result (section + metric)
% that justifies it.
%
%   H1 — Regularity: choreographies with stable temporal rhythms → better outcomes.
%     TODO: Relate Steady_Worker's high active_days and consistent time-of-day
%     distribution to its higher mean final grade (Section~\ref{sec:rq3-results}).
%
%   H2 — Balanced participation: sync/async balance → better outcomes.
%     TODO: Evaluate Balanced_Engager (Forum_View, Synchronous_Join) vs
%     other clusters on grade outcomes.
%
%   H\_Robustness — Context stability: structured subjects → more stable
%   choreographies across cohorts.
%     TODO: Discuss the cosine similarity results and grade-band stickiness
%     from Section~\ref{sec:rq2-results}. Note that the synthetic generator
%     does not vary by subject, so this hypothesis can only be partially assessed.
%
%   H3 — Feedback-centric interaction: frequent feedback → better outcomes.
%     TODO: Evaluate Feedback_Review and related features in the permutation
%     importance results (Section~\ref{sec:rq5-results}).
%
%   H4 — Purposeful navigation: efficient navigation → stronger progress.
%     TODO: Discuss the role of Page_Navigation in Low_Engagement (over-represented
%     there) vs its importance in the outcome model.
%
%   H5 — Risk signals: inactivity / last-minute bursts → weaker outcomes.
%     TODO: Discuss Night_Crammer (Evening/Night over-represented, lower pass
%     rate) and Low_Engagement (low volume, low pass rate) as evidence for H5.

\section{Relation to Prior Work}
\label{sec:relation-to-prior-work}
% TODO: Compare the findings to the literature surveyed in Chapters 2 and 3:
%   - Do the discovered choreographies align with behavioral types described
%     in prior EDM clustering studies?
%   - Does the ARI/NMI validation compare favorably to similar synthetic
%     validation studies in the EDM literature?
%   - How do the outcome associations (RQ3) relate to Beck's findings on the
%     role of at-risk status and parental engagement?
%   - Does the early prediction result (RQ4) echo findings in EDM about the
%     predictive power of first-week engagement signals?
%   - How does the interpretability approach (RQ5) advance on or complement
%     prior feature importance analyses in educational prediction models?

\section{Limitations}
\label{sec:limitations}
% TODO: Discuss each limitation honestly and in enough depth to satisfy a
% critical reader at the thesis defense:
%
%   (1) Synthetic data — the generator produces idealized, well-separated
%       archetypes. Real Pearson logs will have more noise, behaviorally
%       ambiguous students, longitudinal drift, and subject-level variation
%       not captured here. The synthetic study establishes pipeline validity;
%       it does not make empirical claims about actual Connections Academy students.
%
%   (2) Fixed k = 4 — the number of choreographies in real data may differ.
%       The silhouette criterion on synthetic data suggests k = 4, but the
%       real data may support a different number or reveal that the space is
%       continuous rather than discretely clustered.
%
%   (3) Feature coverage — the act__ and time__ blocks capture what actions
%       are taken and when, but not interaction quality (e.g., length or
%       coherence of forum posts, quality of submitted work, response latency
%       to teacher feedback). Richer features may yield better-separated clusters.
%
%   (4) At-risk confounding — the generator correlates at-risk status with
%       archetype (Low_Engagement is predominantly at-risk). This structural
%       correlation limits the independence of the RQ3 analysis; it is
%       difficult to attribute the grade gap purely to choreography when
%       at-risk status and choreography are partially confounded by design.
%
%   (5) Single-semester snapshot — the dataset covers one academic year of
%       synthetic data for one cohort. Longitudinal stability across years,
%       or the effect of school transitions (e.g., elementary to middle school),
%       cannot be assessed.

\section{Threats to Validity}
\label{sec:threats-to-validity}
% TODO: Use the standard four-category framework (Wohlin et al.) to structure
% the threats:
%
%   Internal validity:
%     - k-means is sensitive to centroid initialization; fixed seed (42) reduces
%       but does not eliminate this sensitivity. Report whether different seeds
%       produced materially different cluster assignments (if tested).
%     - The dominant-cluster heuristic (mode of weekly labels per student)
%       collapses within-student variability into a single label for RQ3/RQ4.
%       Students who transition between choreographies mid-year are misrepresented.
%
%   External validity:
%     - Synthetic data; results cannot be directly generalized to real
%       Connections Academy students without replication on the real dataset.
%     - The generator was calibrated from the K-12 online learning literature
%       but not from actual Pearson log statistics; the feature distributions
%       may not match the real data's true marginals.
%
%   Construct validity:
%     - "Pass" as a binary outcome label conflates different failure modes
%       (dropout, academic failure, incomplete). Finer-grained labels might
%       reveal different choreography effects.
%     - ARI as a stability metric assumes that cluster label permutations are
%       the only source of disagreement; it may be inflated if the number of
%       students present in both consecutive weeks is small.
%
%   Statistical conclusion validity:
%     - Five RQs tested on the same dataset increase the risk of Type I error;
%       no Bonferroni correction was applied.
%     - ANOVA assumes homogeneity of variance across groups; this was not
%       formally tested (consider adding Levene's test in the final version).
